%% ============================================================
%%  Enescu (2020) — Revisión detallada + aporte al proyecto
%%  Datos extraídos del PDF en Zotero (38 pp, Energies MDPI)
%%  Compilar: pdflatex enescu_review_beamer.tex  (x2)
%% ============================================================
\documentclass[12pt,aspectratio=169]{beamer}

\usepackage[utf8]{inputenc}
\usepackage[T1]{fontenc}
\usepackage[spanish]{babel}
\usepackage{booktabs}
\usepackage{hyperref}
\usepackage{tikz}
\usetikzlibrary{positioning,arrows.meta,shapes.geometric}

%% ── Tema ────────────────────────────────────────────────────
\usetheme{Madrid}
\usecolortheme{whale}
\setbeamertemplate{navigation symbols}{}
\setbeamertemplate{footline}[frame number]

%% Colores personalizados
\definecolor{azulTesis}{RGB}{26,35,126}
\definecolor{verdeDTR}{RGB}{13,71,161}
\definecolor{rojoAccent}{RGB}{183,28,28}
\definecolor{grisClaro}{RGB}{245,245,245}
\definecolor{azulClaro}{RGB}{232,234,246}

\setbeamercolor{title}{bg=azulTesis, fg=white}
\setbeamercolor{frametitle}{bg=azulTesis, fg=white}
\setbeamercolor{block title}{bg=azulTesis!85, fg=white}
\setbeamercolor{block title alerted}{bg=rojoAccent, fg=white}
\setbeamercolor{block title example}{bg=verdeDTR, fg=white}
\setbeamercolor{item projected}{bg=azulTesis, fg=white}

%% ── Metadatos ───────────────────────────────────────────────
\title[Enescu et al.\ 2020]
      {Thermal Assessment of Power Cables and Impacts\\
       on Cable Current Rating: An Overview}
\subtitle{Revisión crítica del paper y aportes al proyecto de tesis MIA}
\author{Enescu, D. \and Colella, P. \and Russo, A.}
\institute[Valahia / Polito]{%
  $^1$University Valahia of Targoviste, Romania\\
  $^2$Politecnico di Torino, Italy}
\date{Energies 2020, 13(20), 5319 \quad|\quad doi:10.3390/en13205319}

%% ============================================================
\begin{document}

%% ── Portada ─────────────────────────────────────────────────
\begin{frame}
  \titlepage
\end{frame}

%% ── Índice ───────────────────────────────────────────────────
\begin{frame}{Contenido}
  \tableofcontents
\end{frame}

%% ============================================================
\section{Identificación del paper}

\begin{frame}{Datos bibliográficos completos}
  \begin{columns}[T]
    \column{0.52\textwidth}
      \begin{block}{Referencia}
        \small
        \textbf{Enescu, D., Colella, P., \& Russo, A.} (2020).\\
        \textit{Thermal assessment of power cables and impacts on cable
        current rating: An overview.}\\
        \textbf{Energies}, 13(20), 5319.\\[0.3em]
        Recibido: 14 ago 2020 \quad Aceptado: 10 oct 2020\\
        Publicado: 13 oct 2020\\[0.3em]
        \href{https://doi.org/10.3390/en13205319}
             {\texttt{doi:10.3390/en13205319}}
      \end{block}
      \begin{exampleblock}{Tipo de paper}
        \small Revisión sistemática (\textit{overview}) --- 38 páginas\\
        Revista: \textit{Energies} (MDPI) --- Open Access\\
        113 referencias bibliográficas
      \end{exampleblock}
    \column{0.44\textwidth}
      \begin{block}{Keywords}
        \small
        \begin{itemize}
          \item cable rating / ampacity
          \item electric cables
          \item heat transfer
          \item thermal model
          \item finite difference method (FDM)
          \item finite element method (FEM)
          \item probabilistic models
          \item harmonic distortion
        \end{itemize}
      \end{block}
  \end{columns}
\end{frame}

%% ============================================================
\section{Autores e instituciones}

\begin{frame}{Autores, afiliaciones y financiamiento}
  \begin{columns}[T]
    \column{0.54\textwidth}
      \begin{block}{Autores y afiliaciones}
        \small
        \textbf{Diana Enescu$^{1,*}$}\\
        Dept. of Electronics, Telecommunications and Energy\\
        \textit{University Valahia of Targoviste}\\
        Targoviste 130004, \textbf{Romania}\\
        \href{mailto:diana.enescu@valahia.ro}
             {\texttt{diana.enescu@valahia.ro}} (correspondencia)\\[0.6em]
        \textbf{Pietro Colella$^{2}$} \& \textbf{Angela Russo$^{2}$}\\
        Dipartimento Energia ``Galileo Ferraris''\\
        \textit{Politecnico di Torino}\\
        Torino 10129, \textbf{Italy}
      \end{block}
    \column{0.42\textwidth}
      \begin{alertblock}{Financiamiento}
        \small
        ``\textit{This research received \textbf{no external funding.}}''\\[0.5em]
        Trabajo académico independiente sin sponsor industrial ni agencia
        gubernamental.
      \end{alertblock}
      \vspace{0.4em}
      \begin{block}{Contribuciones de autores}
        \small
        Todos los autores contribuyeron por igual en:\\
        conceptualización, metodología, análisis formal, investigación,
        redacción, revisión y visualización.
      \end{block}
  \end{columns}
\end{frame}

%% ============================================================
\section{Objetivos del paper}

\begin{frame}{Objetivos declarados}
  \begin{block}{Objetivo principal (p.\,2)}
    \small
    \textit{``Revisar los modelos térmicos de cables subterráneos
    indicando las evoluciones ocurridas, desde modelos básicos con hipótesis
    generales hacia especificaciones más detalladas para casos prácticos.''}
  \end{block}
  \vspace{0.4em}
  \begin{columns}[T]
    \column{0.48\textwidth}
      \textbf{Objetivos específicos:}
      \begin{enumerate}\small
        \item Resumir conceptos de transferencia de calor para el análisis
              térmico del cable
        \item Presentar la analogía electrotérmica y los modelos equivalentes
              de circuito
        \item Revisar métodos analíticos y numéricos (FDM, FEM, COMSOL, ANSYS)
        \item Relacionar modelos térmicos con el cálculo de la corriente
              admisible (ampacidad)
        \item Analizar efectos avanzados: distorsión armónica y modelos
              probabilísticos
      \end{enumerate}
    \column{0.48\textwidth}
      \textbf{Justificación (gap identificado):}
      \begin{alertblock}{}
        \small
        No existía una revisión reciente que tratara \textbf{simultáneamente}
        modelado del suelo, modelos térmicos, cálculo de ampacidad,
        efectos eléctricos avanzados y pronóstico de corriente admisible.
      \end{alertblock}
  \end{columns}
\end{frame}

\begin{frame}{Enfoque metodológico del paper}
  \vspace{-0.2em}
  \begin{columns}[c]
    %% ── Problema ──────────────────────────────────────────
    \column{0.18\textwidth}
      \centering
      \colorbox{rojoAccent!18}{\parbox{0.95\linewidth}{%
        \centering\small\color{rojoAccent}\textbf{PROBLEMA}\\[3pt]
        \footnotesize\color{black}
        ¿Cuál es la\\máx.\ corriente\\admisible del\\cable?
      }}

    %% ── Flecha ────────────────────────────────────────────
    \column{0.04\textwidth}
      \centering\large$\Rightarrow$

    %% ── Bloques metodológicos ──────────────────────────────
    \column{0.50\textwidth}
      \setlength{\fboxsep}{4pt}
      \fcolorbox{azulTesis}{azulClaro}{\parbox{\linewidth}{%
        \footnotesize\textbf{Sec.\,2 --- Transferencia de calor}\\[-1pt]
        \tiny Ecuación de Fourier $\cdot$ Conducción, convección, radiación $\cdot$
        Propiedades del suelo ($k$, humedad)
      }}\\[4pt]
      \fcolorbox{azulTesis}{azulClaro}{\parbox{\linewidth}{%
        \footnotesize\textbf{Sec.\,3 --- Modelos térmicos}\\[-1pt]
        \tiny Analogía electrotérmica ($R_t$, $C_t$) $\cdot$ Métodos analíticos
        $\cdot$ FDM y FEM (COMSOL, ANSYS)
      }}\\[4pt]
      \fcolorbox{azulTesis}{azulClaro}{\parbox{\linewidth}{%
        \footnotesize\textbf{Sec.\,4 --- Corriente admisible}\\[-1pt]
        \tiny IEC~60287 (permanente) $\cdot$ IEC~60853 (transitorio) $\cdot$
        Armónicos $\cdot$ Modelos probabilísticos
      }}

    %% ── Flecha ────────────────────────────────────────────
    \column{0.04\textwidth}
      \centering\large$\Rightarrow$

    %% ── Resultado ─────────────────────────────────────────
    \column{0.20\textwidth}
      \centering
      \colorbox{verdeDTR!20}{\parbox{0.95\linewidth}{%
        \centering\small\color{verdeDTR}\textbf{RESULTADO}\\[3pt]
        \footnotesize\color{black}
        Estado del\\arte $+$\\recomend.\\DTR
      }}
  \end{columns}

  \vspace{0.5em}
  \begin{alertblock}{}
    \centering\small
    Estrategia: fundamentos físicos $\to$ modelos térmicos $\to$
    cálculo normativo, integrando métodos analíticos y numéricos.
  \end{alertblock}
\end{frame}

%% ============================================================
\section{Contenido y estructura}

\begin{frame}{Estructura del paper (38 pp, 5 secciones)}
  \centering
  \begin{tikzpicture}[>=Stealth, node distance=0.55cm,
      sec/.style={draw=azulTesis, rounded corners=3pt, fill=azulClaro,
                  text width=3.5cm, align=center, font=\footnotesize,
                  inner sep=5pt, minimum height=0.85cm},
      sub/.style={draw=gray!50, rounded corners=2pt, fill=grisClaro,
                  text width=3.2cm, align=left, font=\tiny,
                  inner sep=4pt}]

    \node[sec] (s1) {Sec.\,1\\Introducción};
    \node[sec, right=0.4cm of s1] (s2) {Sec.\,2\\Transferencia de Calor};
    \node[sec, right=0.4cm of s2] (s3) {Sec.\,3\\Modelos Térmicos};
    \node[sec, right=0.4cm of s3] (s4) {Sec.\,4\\Corriente Admisible};
    \node[sec, right=0.4cm of s4] (s5) {Sec.\,5\\Conclusiones};

    \node[sub, below=0.4cm of s2]
      {Cond.\ Fourier, conv., rad.\\Propiedades del suelo\\Saturación/humedad};
    \node[sub, below=0.4cm of s3]
      {Analogía electrotérmica\\Métodos: FDM, FEM\\COMSOL, ANSYS};
    \node[sub, below=0.4cm of s4]
      {IEC 60287 (steady-state)\\IEC 60853 (transitorio)\\Armónicos, probabilidad};
    \node[sub, below=0.4cm of s5]
      {FEM 3D viable\\Modelos simpl.\ emergentes\\DTR como aplicación clave};

    \draw[->, azulTesis] (s1) -- (s2);
    \draw[->, azulTesis] (s2) -- (s3);
    \draw[->, azulTesis] (s3) -- (s4);
    \draw[->, azulTesis] (s4) -- (s5);
  \end{tikzpicture}
\end{frame}

\begin{frame}{Principales hallazgos del paper}
  \begin{columns}[T]
    \column{0.50\textwidth}
      \begin{block}{Modelos térmicos del cable}
        \begin{itemize}\small
          \item Ecuación de Fourier 2D con coeficientes variables $k(x,y)$
          \item Analogía electrotérmica: $R_t$ (resistencias), $C_t$ (capacitancias)
          \item Modelos de red térmica (ladder network)
          \item FEM 2D/3D con COMSOL y ANSYS: alta precisión pero costoso
        \end{itemize}
      \end{block}
      \begin{block}{Suelo y backfill}
        \begin{itemize}\small
          \item $k_\text{soil}$ varía con humedad, densidad y mineralogía
          \item Zona seca alrededor del cable (drying-out): reduce ampacidad
          \item Backfill especial: mejora conductividad térmica
        \end{itemize}
      \end{block}
    \column{0.46\textwidth}
      \begin{block}{Corriente admisible (ampacidad)}
        \begin{itemize}\small
          \item \textbf{IEC 60287}: régimen estacionario, $k$ homogénea, método
                analítico estándar
          \item \textbf{IEC 60853}: régimen transitorio/cíclico
          \item FEM supera a IEC cuando la geometría es compleja
          \item Armónicos elevan $R_\text{AC}$ y reducen ampacidad
        \end{itemize}
      \end{block}
      \begin{alertblock}{Tendencia emergente (conclusiones)}
        \small
        DTR (\textit{dynamic thermal rating}) reconocido como la aplicación
        que permite mejor explotación de los cables en tiempo real.
      \end{alertblock}
  \end{columns}
\end{frame}

%% ============================================================
\section{Definiciones del paper}

\begin{frame}{Definiciones fundamentales (Sec.\,1--2)}
  \begin{block}{Corriente admisible / Ampacidad (p.\,2)}
    \small
    \textit{``The current rating (also denoted as \textbf{ampacity}) refers to the
    maximum current at which the cable can operate without exceeding the
    temperature limits for the insulation material.''}
  \end{block}
  \vspace{0.3em}
  \begin{block}{Clasificación térmica dinámica --- DTR (p.\,2)}
    \small
    \textit{``[Dynamic cable rating] consists of the determination of the
    \textbf{maximum permissible current loading} of the cable during the time.
    Dynamic rating involves an iterative method for continuous calculation of
    the conductor temperature by considering the history of the current loading,
    alongside the parameters and thermal conditions of the cable and the
    environment.''}
  \end{block}
  \vspace{0.3em}
  \begin{exampleblock}{Mecanismos de transferencia de calor (p.\,4)}
    \small
    \begin{itemize}
      \item \textbf{Conducción}: transferencia por contacto directo; rige en el interior del cable y en el suelo
      \item \textbf{Convección}: relevante en cables en aire o agua; no aplica en cables enterrados
      \item \textbf{Radiación}: transferencia por ondas EM; solo en cables en aire (\textit{``does not necessitate a medium''})
    \end{itemize}
  \end{exampleblock}
\end{frame}

\begin{frame}{Definiciones de modelos y magnitudes (Sec.\,2--3)}
  \begin{columns}[T]
    \column{0.50\textwidth}
      \begin{block}{Fuentes de calor internas $\dot{Q}_\text{gen}$ (p.\,4)}
        \small
        \textit{``Heat flow rate generated inside a specific power cable,
        by \textbf{Joule, dielectric and ferromagnetic losses}.''}\\[0.4em]
        Pérdidas en el conductor: efecto Joule\\
        Pérdidas en la aislación: pérdidas dieléctricas\\
        Pérdidas en pantallas metálicas: corrientes inducidas
      \end{block}
      \vspace{0.4em}
      \begin{block}{Estabilidad térmica del suelo (p.\,12)}
        \small
        \textit{``The ability of the soil to keep its \textbf{thermal resistivity
        constant} in the presence of a heat source is named \textbf{thermal
        stability} of the soil.''}\\[0.3em]
        Su degradación genera la zona seca (\textit{drying-out}) que
        reduce la ampacidad del cable.
      \end{block}
    \column{0.46\textwidth}
      \begin{block}{Analogía electrotérmica (Sec.\,3)}
        \small
        Mapa de equivalencia calor $\leftrightarrow$ circuito eléctrico:
        \begin{center}\small
        \begin{tabular}{ll}
          \toprule
          \textbf{Magnitud térmica} & \textbf{Análoga eléctrica}\\
          \midrule
          Flujo de calor $\dot{Q}$ & Corriente $I$\\
          Temperatura $T$ & Tensión $V$\\
          Resist. térmica $R_t$ & Resistencia $R$\\
          Capacit. térmica $C_t$ & Capacitancia $C$\\
          \bottomrule
        \end{tabular}
        \end{center}
        Permite resolver la EDP de calor como un circuito RLC.
      \end{block}
      \begin{exampleblock}{Ecuación de Fourier (régimen permanente)}
        \small
        $\nabla \cdot (k\,\nabla T) + \dot{q}_\text{gen} = 0$\\[0.2em]
        con $k$ constante $\Rightarrow$ Laplaciano: $k\,\nabla^2 T + \dot{q} = 0$
      \end{exampleblock}
  \end{columns}
\end{frame}

%% ============================================================
\section{Aporte al proyecto de tesis}

\begin{frame}{Posicionamiento de la tesis respecto a Enescu 2020}
  \begin{block}{Aporte directo al proyecto}
    \centering
    \textit{``Establece el estado del arte que la tesis amplía,
    posicionando el enfoque PINN como alternativa de mayor precisión
    frente a los modelos analíticos (IEC 60287) y de elementos finitos
    existentes, sin incremento del costo computacional.''}
  \end{block}
  \vspace{0.5em}
  \begin{columns}[T]
    \column{0.48\textwidth}
      \textbf{Lo que Enescu 2020 deja abierto:}
      \begin{itemize}\small
        \item No considera $k_\text{soil}(x,y)$ \textbf{espacialmente variable}
              en el mismo solver
        \item FEM requiere remallado y solver iterativo por cada escenario
        \item No integra inferencia de parámetros (problema inverso)
        \item Sin capacidad de evaluación en \textit{tiempo real}
      \end{itemize}
    \column{0.48\textwidth}
      \textbf{Respuesta de la tesis (PINN):}
      \begin{itemize}\small
        \item Resuelve la EDP con $k(x,y)$ heterogéneo (backfill PAC + suelo)
              en \textbf{un solo entrenamiento}
        \item Inferencia $<$\,1\,ms (vs.\ FEM: minutos por simulación)
        \item Error PINN\,vs.\,FEM $<$\,1\,K en benchmarks validados
        \item Marco extensible a problemas inversos ($k$ desconocida)
      \end{itemize}
  \end{columns}
\end{frame}

\begin{frame}{Tabla de posicionamiento}
  \centering\small
  \begin{tabular}{lccc}
    \toprule
    \textbf{Aspecto} & \textbf{IEC 60287} & \textbf{FEM} & \textbf{PINN (tesis)} \\
    \midrule
    Régimen        & Permanente     & Ambos          & Permanente (ext.\ trans.) \\
    $k$ del suelo  & Homogénea      & Variable       & $k(x,y)$ heterogénea \\
    Geometría      & Simplificada   & Arbitraria     & Arbitraria \\
    Velocidad      & Instantánea    & Lenta (min)    & \textbf{Rápida} ($<$1\,ms) \\
    Precisión      & $\pm$5--10\,K  & $\pm$0.5\,K    & $\pm$0.6--1\,K \\
    Prob.\ inverso & No             & No             & Sí (potencial) \\
    \bottomrule
  \end{tabular}
  \vspace{0.5em}
  \begin{alertblock}{}
    \centering
    Enescu 2020 valida que \textbf{FEM es la referencia de precisión}:\\
    la tesis demuestra que el PINN alcanza esa precisión
    con la velocidad de un modelo analítico.
  \end{alertblock}
\end{frame}

%% ── Cierre ──────────────────────────────────────────────────
\begin{frame}{Referencia}
  \begin{thebibliography}{1}\small
    \bibitem{enescu2020}
      D.~Enescu, P.~Colella, A.~Russo (2020).
      \textit{Thermal assessment of power cables and impacts on cable
      current rating: An overview.}
      \textbf{Energies}, 13(20), 5319.\\
      Recibido: 14-ago-2020 \quad Publicado: 13-oct-2020\\
      \href{https://doi.org/10.3390/en13205319}
           {doi:10.3390/en13205319}\\[0.4em]
      \footnotesize
      $^1$ University Valahia of Targoviste, Romania\\
      $^2$ Politecnico di Torino, Italy\\
      Financiamiento: ninguno (sin fuente externa)
  \end{thebibliography}
\end{frame}

\end{document}
